\documentclass[a4paper, 11pt]{article}

% --- Pacchetti Fondamentali ---
\usepackage[utf8]{inputenc}
\usepackage[T1]{fontenc}
\usepackage{ebgaramond} % Font elegante per il testo
\usepackage[italian]{babel}
\usepackage{gillius2} % Font moderno per titoli e commenti
\usepackage{geometry}
\geometry{margin=2.5cm}
\usepackage[svgnames]{xcolor}
\usepackage{multicol, verse, changepage}
\usepackage{enumitem}
\usepackage{xparse}
\usepackage{hyperref}
\usepackage{caption}

% --- Configurazione Numerazione e Indice ---
\setcounter{secnumdepth}{4}
\setcounter{tocdepth}{4}

\hypersetup{
    colorlinks=true,
    linkcolor=MediumBlue,
    urlcolor=DarkCyan,
    pdftitle={Letteratura Italiana A - Edoardo Lo Iacono}
}

% --- Comandi Personalizzati ---

% Titolo della poesia pulito e moderno
\newcommand{\mypoemtitle}[2]{%
    \vspace{2em}
    \begin{flushleft}
        {\sffamily\color{Tomato!80}\scriptsize\bfseries\MakeUppercase{#2}} \\ 
        \vspace{2pt}
        {\sffamily\LARGE\bfseries #1} \\ 
        \vspace{-8pt}
        {\color{Tomato!40}\rule{\textwidth}{0.8pt}} 
    \end{flushleft}
    \vspace{1em}
}

% Ambiente per i commenti a due colonne
\makeatletter
\newcommand\thickhrulefill[1][0.5pt]{\leavevmode\leaders\hrule height #1\hfill\kern\z@}
\makeatother

\newenvironment{commentaries}{%
    \nopagebreak
    {\color{Gray!30}\thickhrulefill}\vspace{-1ex}
    \sffamily\small\begin{multicols*}{2}
    \description[leftmargin=1.5em, style=unboxed, font=\sffamily\bfseries\color{Tomato}]
}{%
    \enddescription\end{multicols*}
}

% Comando per citare i versi (scolio)
% Modifica del comando \scolio per accettare testo libero nelle []
\NewDocumentCommand{\scolio}{O{}m}{\item[#1] \textbf{#2}:}

\title{Letteratura Italiana A}
\author{A cura di: Edoardo Lo Iacono}
\date{}

\begin{document}

\maketitle

\vfill 

\begin{center}
    A.S. 2025 - 2026 \\
    Lezioni del prof. \href{https://cdslettere.campusnet.unito.it/do/docenti.pl/Alias?valter.boggione#tab-profilo}{Valter Boggione} + Integrazione dei testi
\end{center}

\newpage
\tableofcontents
\newpage

\section{La scuola siciliana}
\subsection{Le origini della letteratura in Italia}
La letteratura si sviluppa in Italia molto tardi rispetto ad altri paesi come la Francia o la Germania, nonostante questo sviluppo tardivo nasce però, tramite un processo di importazione, una letteratura già matura.
Per comprendere le motivazioni alla base di questo rallentamente è importante considerare due elementi fondamentali:
\begin{itemize}
    \item Quando la letteratura si origina, l'Italia è ancora fortemente frammentata, è quindi difficile stabilire una letteratura Italiana
    \item La presenza del vaticano e quindi il forte potere culturale esercitato dalla chiesa favoriscono una cultura classica e l'utilizzo della lingua latina
\end{itemize}
Possiamo dire che a determinare l'inizio della letteratura in Italia grazie ad un fenomeno di importazione sia stato un preciso evento storico: la \textbf{crociata contro gli Albigesi}
Gli Albigesi, o catari erano una popolazione del sud della Francia considerata eretica in quanto non apprezzava l'uso dei beni materiali perpetuato dalla chiesa Romana. La crociata ai loro danni durò circa 20 anni: dal 1200 al 1220 e portò molti Albigesi a migrare in altre parti d'Euorpa in cerca di protezione.
Questa crociata era portata avanti dalla collaborazione tra :
\begin{itemize}
    \item \textit{Innocenzo III}: Papa che voleva scardinare i catari in quanto eretici
    \item \textit{Luigi IX}: Re di Francia che voleva unificare i territori del sud in quel momento particolarmente frammentati. E' importante però, anche in quest'ultimo il valore religioso, Luigi IX era infatti soprannominato \textit{"il santo"}
\end{itemize}
A causa delle persecuzioni sul piano religioso, i Catari si trovarono a cercare rifugio e protezione in cui erano protetti da questo punto di vista, andarono quindi in Germania ed in molti altri luoghi tra cui l'Italia del Sud in cui prosperava la corte di \textbf{Federico II di Svevia}, una corte che presentava delle importantissime caratteristiche favorevoli ai Catari ovvero una forte \textit{apertura culturale} ed il fatto di essere una \textit{corte laica}.
Questa stessa corte così aperta dal punto di vista ideolgico si rivela però fortemente assolutista dal punto di vista politico e questo ha varie ripercussioni nel modo di fare poesia.
\subsection{Caratteristiche e differenze con la poesia provenzale}
Con l'avvento della Scuola Siciliana cambiano infatti molti paradigmi della precedente poesia provenzale:
In primo luogo i poeti non fanno più quello di professione ma sono invece \textit{funzionari di Federico II}, questo comporta un cambiamento nella concezione della poesia stessa, la quale viene ora approcciata come un passatempo individuale e non più destinato alle grandi rappresentazione. Questo comporta inoltre un scissione tra la figura del \textbf{cantante} e quella del \textbf{poeta}, scissione che non muterà fino ai moderni cantautori. Questo cambiamento, per quanto insignificante possa apparire apre in realtà le porte ad una nuova concezione della metrica che accompagnerà tutta la poesia Italiana delle origni; viene infatti applicata una metrica molto \textbf{più rigida} in quanto non si può con la voce andare ad aggiustare l'accento o la lunghezza dei versi, abbiamo così un'irrigidimento della metrica e una selezione di schemi metrici. In particolare nella scuola Siciliana troviamo solamente \textbf{canzoni}, alcune sporadiche \textbf{ballate} ed il metro da essi inventato: \textbf{il sonetto}. Anche le tematiche, proprio come la metrica, si vanno ad irrigidire; infatti visto l'assolutismo della scuola Siciliana e visto il ruolo di funzionari esercitato dai "poeti" diviene totalmente vietato parlare di politica, cadono così interi generi come quello del \textit{sirventese} (presente ancora in altre parti d'Italia in particolare quella settentrionale), la poesia ricopre un ruolo maggiormente di svago concentrandosì così sui piaceri edonistici ed in particolar modo sull'\textbf{amore}, quest'ultimo addirittura paragonato e posto innanzi ai valori della chiesa, decisione che non porta alcuna ripercussione in questi poeti visto il contesto fortemente laico. 
\subsection{Gli scrittori ed il siciliano illustre}
La scuola siciliana viene definita \textbf{scuola} in quanto rappresenta un gruppo di intellettuali che condividono le stesse idee avvalendosi come abbiamo detto di uno stile e di tematiche molto rigide. Queste premesse portano gli autori della scuola siciliana ad essere sostanzialmente \textit{anonimi} ed irriconoscibili tra di loro, tranne nei rari casi in cui è presente come vedremo una firma. L'attribuzione dei testi non è ancora oggi completamente sicura, è inoltre fondamentale notare come nel leggere i loro testi non ci troviamo più davanti alla versione originale ma ad una versione \textbf{toscanizzata} dai copisti toscani, modifica che renderà i testi più accessibili portando anche alla corrente Siculo - Toscana, andando però a perdere completamente il suo valore lessicale; è infatti impossibile applicare una traduzione inversa in quanto le poesie dei siciliani erano scritte in un \textbf{siciliano illustre}, ovvero una lingua colta usata solamente in ambito letterario che univa in se elemente delle più importanti lingue della letteratura del tempo: il latino, il provenzale ed il Siciliano.
La toscanizzazione dei componimenti è palese data la presenza di rime improprie chiamate \textbf{rime siciliane}.
I principali autori di questa corrente si caratterizzano per delle differenze minime e sono:
\subsubsection{Giacomo da Lentini}
Giacomo da Lentini è il poeta più importante della scuola siciliana, considerato da molti come fondatore di questa stessa corrente, va a lui attribuita l'invenzione del \textbf{sonetto}.
Per quanto riguarda la sua vita non si hanno molte informazioni biografiche, sappiamo però per fonte diretta e tramite Dante che faceva di professione il Notaio, è inoltre l'autore più prolifico ma come detto non c'è la certezza che tutti i testi a lui attribuiti siano esattamente dell'autore.
Tra i testi attribuiti a Giacomo da Lentini troviamo:
\begin{itemize}
    \item \hyperlink{para:meravigliosamente}{Meravigliosamente}
    \item \hyperlink{para:postoincore}{Io m'aggio posto in core a Dio servire}
    \item \hyperlink{para:amorun}{Amor è un disio che ven da core}
\end{itemize}
\subsubsection{Stefano Protonotaro}
Anche di Stefano Protonotaro abbiamo poche informazioni biografiche, lo studiamo per un motivo particolare, la sua importanza poetica è infatti minore rispetto a quella di Giacomo da Lentini ma è importante in quanto è l'unico poeta di cui ci rimane un testo scritto in \textit{siciliano illustre} e non toscanizzato.
La poesia in questione è: \hyperlink{para:pirmeu}{Pir meu cori alligrari
}
\subsubsection{Cielo d'Alcamo}
Cielo d'Alcamo è uno degli esponenti della scuola siciliana, la suo opera più importante: \hyperlink{par:rosafresca}{Rosca fresca aulentissima} riprende le tematiche tipiche della poesia Siciliana in maniera ironica e scherzosa
\newpage
\hypertarget{para:meravigliosamente}{}

\mypoemtitle{Meravigliosamente - Giacomo da Lentini}{Canzonetta}

\textbf{Commento introduttivo:}\\
In questa canzonetta, Giacomo da Lentini affronta uno dei temi cardinali della poesia Siciliana ovvero l'espressione dell'amore come esperienza visiva, viene applicato inoltre un altro topos letterario tipico della letteratura provenzale ovvero quello dell'amante timido ed incapace di rivelare i propri sentimenti alla donna amata.

\vspace{2em}

\begin{adjustwidth}{2em}{0pt}
\poemlines{5}
\verselinenumbersleft
\begin{verse}
Meravigliosa - mente \label{m1} \\
un amor mi distringe \label{m2} \\
e mi tene ad ogn’ora. \label{m3} \\
Com’om che pone mente \label{m4} \\
in altro exemplo \textbf{pinge} \label{m5} \\
la simile \textbf{pintura}, \label{m6} \\
così, bella, facc’eo, \label{m7} \\
che ’nfra lo core meo \label{m8} \\
porto la tua \textbf{figura}. \label{m9} \\!

In cor par ch’eo vi porti, \label{m10} \\
\textbf{pinta} come parete, \label{m11} \\
e non pare difore. \label{m12} \\
O Deo, co’ mi par forte \label{m13} \\
non so se lo sapete, \label{m14} \\
con’ v’amo di bon core; \label{m15} \\
ch’eo son sì vergognoso \label{m16} \\
ca pur vi guardo ascoso \label{m17} \\
e non vi mostro amore. \label{m18} \\!

Avendo gran disio \label{m19} \\
\textbf{dipinsi una pintura}, \label{m20} \\
bella, voi simigliante, \label{m21} \\
e quando voi non vio \label{m22} \\
guardo ’n quella \textbf{figura}, \label{m23} \\
par ch’eo v’aggia davante: \label{m24} \\
come quello che crede \label{m25} \\
salvarsi per sua fede, \label{m26} \\
ancor non veggia inante. \label{m27} \\!

Al cor m’ard’una doglia, \label{m28} \\
com’ om che ten lo foco \label{m29} \\
a lo suo seno ascoso, \label{m30} \\
e quando più lo ’nvoglia, \label{m31} \\
allora arde più loco \label{m32} \\
e non pò star incluso: \label{m33} \\
similemente eo ardo \label{m34} \\
quando pass’e non guardo \label{m35} \\
a voi, vis’amoroso. \label{m36} \\!

S’eo guardo, quando passo, \label{m37} \\
inver’ voi no mi giro, \label{m38} \\
bella, per risguardare; \label{m39} \\
andando, ad ogni passo \label{m40} \\
getto uno gran sospiro \label{m41} \\
ca facemi ancosciare; \label{m42} \\
e certo bene ancoscio, \label{m43} \\
c’a pena mi conoscio, \label{m44} \\
tanto bella mi pare. \label{m45} \\!

Assai v’aggio laudato, \label{m46} \\
madonna, in tutte parti, \label{m47} \\
di bellezze c’avete. \label{m48} \\
Non so se v’è contato \label{m49} \\
ch’eo lo faccia per arti, \label{m50} \\
che voi pur v’ascondete: \label{m51} \\
sacciatelo per singa \label{m52} \\
zo ch’eo no dico a linga, \label{m53} \\
quando voi mi vedite \label{m54} \\!

Canzonetta novella, \label{m55} \\
va’ canta nova cosa; \label{m56} \\
lèvati da maitino \label{m57} \\
davanti a la più bella, \label{m58} \\
fiore d’ogn’amorosa, \label{m59} \\
bionda più c’auro fino: \label{m60} \\
«Lo vostro amor, ch’è caro, \label{m61} \\
donatelo al Notaro \label{m62} \\
ch’è nato da Lentino». \label{m63}
\end{verse}
\end{adjustwidth}

\vspace{3em}

\begin{commentaries}
    \scolio[\ref{m1}]{Meravigliosamente} L'avverbio definisce lo stupore dell'amante davanti alla bellezza.
    \scolio[m4]{Com'omo che pone mente} Similitudine col pittore che osserva un modello per ritrarlo.
    \scolio[m8]{'nfra lo core meo} L'immagine della donna custodita interiormente.
    \scolio[m11]{pinta come parete} Come se l'immagine fosse affrescata nel cuore.
    \scolio[m25]{come quello che crede} Similitudine religiosa: l'amore come fede in ciò che non si vede direttamente.
    \scolio[m62]{al Notaro} Firma del poeta, identificato come Giacomo da Lentini.
\end{commentaries}
\newpage
\hypertarget{para:postoincore}{}

\mypoemtitle{Io m’aggio posto in core a Dio servire - Giacomo da Lentini}{Sonetto}

\textbf{Commento introduttivo:}\\
In questo sonetto, Giacomo da Lentini affronta il rapporto tra amore profano e beatitudine celeste, teorizzando la necessità della presenza della donna anche nel regno di Dio per raggiungere la vera gioia.

\vspace{1.5em}

\begin{adjustwidth}{2em}{0pt}
\poemlines{5}
\verselinenumbersleft
\begin{verse}
Io m’aggio posto in core a Dio servire, \label{posto1}\\
com’io potesse gire in paradiso, \label{posto2}\\
al santo loco, c’aggio audito dire, \\
o’ si mantien sollazzo, gioco e riso. \label{posto4}\\!

     Sanza mia donna non vi voria gire, \\
quella c’à blonda testa e claro viso, \label{posto6}\\
che sanza lei non poteria gaudere, \\ 
estando da la mia donna diviso. \\!

     Ma no lo dico a tale intendimento, \\
perch’io pecato ci volesse fare; \\
se non veder lo suo bel portamento \\!

     e lo bel viso e ’l morbido sguardare: \\
che·l mi teria in gran consolamento, \\
veggendo la mia donna in ghiora stare. \label{posto14}\\
\end{verse}
\end{adjustwidth}

\vspace{1.5em}

\begin{commentaries}
    \scolio[\ref{posto2}]{gire} Andare.
    \scolio[\ref{posto6}]{blonda testa...} Canone di bellezza che diverrà standard nella lirica.
    \scolio[\ref{posto14}]{In ghiora stare} Stare nella gloria del paradiso. ghiora è un termine toscano sicuramente inserito dai copisti
\end{commentaries}

\vspace{0.5em}
{\color{Tomato!40}\hrule height 0.5pt}
\vspace{1em}

\subsubsection*{Analisi}
\begin{quote}
\small \itshape
A partire dalla prima quartina del sonetto, il poeta mette in luce la sua intenzione ovvero quella di servire ed avere fede in Dio per poter andare in paradiso e scoprire se è vero quello che gli è stato detto, ovvero che il paradiso sia un luogo felice e gioioso. già nella prima quartina, ed in particolar modo al vv.4, si crea una relazione tra la chiesa e la corte andando a vedere il paradiso come un ambiente di corte quasi giullaresco.

Nella seconda quartina viene introdotto un'ulteriore elemento presente nella poesia Sicliana, ovvero l'amore messo sullo stesso livello, se non su un livello superiore rispetto a quello della chiesa: Giacomo dice infatti di voler andare in paradiso solo a condizione che ci sia la la sua donna perché essendo da ella diviso non potrebbe essere felice nemmeno in paradiso

Nelle terzine il poeta mette le mani avanti, non vuole la donna per poter fare cose impure ma per le sue qualità spirituali e visive, in modo da poterla contemplare.
\end{quote}
\newpage
\hypertarget{para:amorun}{}

\mypoemtitle{Amor è un disio che ven da core - Giacomo da Lentini}{Sonetto}

\textbf{Commento introduttivo:}\\
In questo sonetto Giacomo da Lentini applica una fenomenologia dell'amore sull'uomo descrivendo quello che accade quando un'uomo si innamora ed andando ad avvalorare l'importanza della vista in questo processo

\vspace{1.5em}

\begin{adjustwidth}{2em}{0pt}
\poemlines{5}
\verselinenumbersleft
\begin{verse}
Amore è uno desi[o] che ven da’ core \\
per abondanza di gran piacimento; \\
e li occhi in prima genera[n] l’amore \\
e lo core li dà nutricamento. \\!

Ben è alcuna fiata om amatore \label{v5}\\
senza vedere so ’namoramento, \\
ma quell’amor che stringe con furore \\ 
da la vista de li occhi ha nas[ci]mento: \\!

ché li occhi rapresenta[n] a lo core \\
d’onni cosa che veden bono e rio  \\
com’è formata natural[e]mente; \label{v11}\\!

e lo cor, che di zo è concepitore, \\
imagina, e [li] piace quel desio: \\
e questo amore regna fra la gente. \label{v14}\\
\end{verse}
\end{adjustwidth}

\vspace{1.5em}

\begin{commentaries}
    \scolio[\ref{v5}]{fiata} Provenzalismo: volta
    \scolio[\ref{v11}]{naturalmente} L'amore è un'esperienza fisica naturale
    \scolio[\ref{v14}]{regna fra la gente} è diffuso in tutto il mondo
\end{commentaries}

\vspace{0.5em}
{\color{Tomato!40}\hrule height 0.5pt}
\vspace{1em}

\subsubsection*{Analisi}
\begin{quote}
\small 
Il componimento si presenta come risposta ad una \textbf{tenzone} cominciata da Jacopo Mostacchi. A questa tenzone prendono parte Jacopo mostacchi, Giacomo da Lentini e Pier delle Vigne i quali cercano di spiegare in maniera filosofica quale sia la natura dell'amore. In particolare fanno riferimento ai concetti aristotelici di \textbf{sostanza} ed \textbf{accidente}. Per Aristotele la sostanza di un oggetto è quella che lo rende tale, non può quindi esistere tale oggetto senza la sua sostanza, un accidente è invece un'attributo di questo elemento che può variare e pur identificandolo non lo rende o meno tale.
Jacopo Mostacchi sostiene che l'amore sia un accidente, Pier delle Vigne sostiene che sia una sostanza senza cui l'uomo non potrebbe vivere e Giacomo da Lentini si pone un po' a metà tra le due visioni.

La struttura di questo sonetto è abbastanza classica, troviamo infatit nella prima quartina la tesi, nella seconda un'antitesi e nelle terzine degli argomenti volti a dare valore alla tesi.

Già a partire dal vv.1 possiamo notare la natura didascalica di questo componimento: Giacomo sembra spiegare al lettore cosa sia l'amore e come funzioni mettendo in luce due componenti fondamentali: \textbf{il cuore} in cui il sentimento risiede e da cui viene nutrito e \textbf{gli occhi} che hanno l'importante compito di generare il sentimento amoroso, tema già affrontato in altri testi dei Siciliani.
Ritornando a considerare la nostra tenzone potremmo sostenere che Giacomo sia più per definire l'amore un'accidente in quanto in assenza della vista, come verrà specificato di seguito, non può esserci vero amore ma l'uomo esiste comunque.

Nella seconda quartina Giacomo sostiene infatti che sia sì possibile innamorarsi senza vedere l'oggetto del proprio amore ma che quell'amore non sia amore vero, critica in questo modo tutta una corrente letteraria e tematica detta \textit{amor de lonh}.

Nelle terzine torna ad avvalorare la sua tesi, dicendo in primo luogo che tutte le cose, belle o brutte, vengono prima filtrate dallo sguardo e per questo motivo l'esperienza amorosa è del tutto naturale, negli ultimi versi sostiene inoltre che sia universale il sentimento amoroso, andando così più in contro alla teoria dell'amore come sostanza
\end{quote}
\newpage
\hypertarget{para:pirmeu}{}

\mypoemtitle{Pir meu cori alligrari - Stefano Protonotaro}{Sonetto}

\textbf{Commento introduttivo:}\\
Questa canzone di Stefano Protonotaro ci permette di vedere la veste originale degli scritti Siciliani, rappresenta infatti un siciliano illustre con al suo interno elementi di diverse lingue letterarie del tempo tra cui latino, provenzale e siciliano, è per questo motivo abbastanza complicata da comprendere dal punto di vista stilistico

\vspace{1.5em}

\begin{adjustwidth}{2em}{0pt}
\poemlines{5}
\verselinenumbersleft
\begin{verse}
Pir meu cori allegrari, \\
chi multu longiamenti\\
senza alligranza e joi d'amuri è statu,\\
mi ritornu in cantari,\\
ca forsi levimenti\\
da dimuranza turniria in usatu\\
di lu troppu taciri;\\
e quandu l'omu à rasuni di diri,\\
ben di' cantari e mustrari alligranza,\\
ca senza dimustranza\\
joi siria sempri di pocu valuri;\\
dunca ben di' cantari onni amaduri.\\
E si per ben amari\\
cantau juiusamenti\\
homo chi avissi in alcun tempu amatu,\\
ben lu diviria fari\\
plui dilittusamenti\\
eu, chi su di tal donna inamuratu,\\
dundi è dulci placiri,\\
preiu e valenza e juiusu pariri\\
e di billizi cuta[n]t' abondanza,\\
chi illu m'è pir simblanza\\
quandu eu la guardu, sintir la dulzuri\\
chi fa la tigra in illu miraturi;\\
chi si vidi livari\\
multu crudilimenti\\
sua nuritura, chi illa à nutricatu,\\
e si bonu li pari\\
mirarsi dulcimenti\\
dintru unu speclu chi li esti amustratu,\\
chi l'ublia siguiri.\\
Cusì m'è dulci mia donna vidiri:\\
chi 'n lei guardandu met[t]u in ublianza\\
tutta'altra mia intindanza,\\
sì chi instanti mi feri sou amuri\\
d'un culpu chi inavanza tutisuri.\\
Di chi eu putia sanar;\\
multu legeramenti,\\
sulu chi fussi a la mia donna a gratu\\
meu sirviri e pinari;\\
m'eu duitu fortimenti\\
chi quandu si rimembra di sou statu\\
nu lli dia displaciri.\\
Ma si quistu putissi adiviniri,\\
ch'Amori la ferissi de la lanza\\
chi mi fer' e mi lanza,\\
ben crederia guarir de mei doluri,\\
ca sintiramu equalimenti arduri.\\
Purriami laudari\\
d'Amori bonamenti,\\
com'omu da lui beni ammiritatu;\\
ma beni è da blasmari\\
Amur virasementi,\\
quandu illu dà favur da l'unu latu,\\
chi si l'amanti nun sa suffiriri,\\
disia d'amari e perdi sua speranza.\\
Ma eu suf[f]ru in usanza,\\
chi ò vistu adessa bon sufrituri\\
vinciri prova et aquistari hunuri.\\
E si pir suffiriri,\\
ni per amar lialmenti e timiri,\\
homo aquistau d'Amur gran beninanza,\\
digu avir confurtanza\\
eu, chi amu e timu e servi[i] a tutt'uri\\
cilatamenti plu chi altru amaduri.\\
\end{verse}
\end{adjustwidth}


\vspace{0.5em}
{\color{Tomato!40}\hrule height 0.5pt}
\vspace{1em}

\subsubsection*{Analisi}
\begin{quote}
\small 
Nonostante la difficoltà linguistica è interessante mettere in evidenza alcuni elementi. In primo luogo la canzone parla ovviamente della gioia prodotta dall'amore.

Al vv.23-24 cita la tigre vista nei bestiari come un animale vanitoso che si perde guardandosi allo specchio, tanche che nelle battute di caccia i cacciatori mettevano sul terreno degli specchi per far stare fermi gli animali.
Ai vv.35-36 si dice invece che la ferita d'amore è brutta solamente se è uno solo a ferirsi e non tutti e due, significa che se l'amore non è corrisposto fa male ma è anche una metafora alla lancia di Achille che, secondo la leggenda, colpendo una volta ferisce e colpendo nuovamente risana
\end{quote}
%Aggiungere rosa fresca aulentissima.

\section{La poesia religiosa}
Abbiamo visto come la letteratura in Italia si sia sviluppata in maniera tardiva e con un modello di importazione francese, questo è vero in particolare al nord ed al sud, possiamo infatti identificare una tripartizione geografica della letteratura in Italia delle orgini, La frammentazione dell'Italia si rispecchia così in quella della letteratura che si divide così:
\begin{itemize}
    \item Al \textbf{nord} troviamo una poesia di stampo francese, sono ancora presenti qui generi di natura politica come il sirventese che sono invece aboliti al sud
    \item Al \textbf{sud} troviamo la scuola siciliana di cui già abbiamo parlato nel capitolo precedente
    \item Al \textbf{centro}, a causa dell'avvento di diversi monasteri si diffonde una poesia fortemente religiosa, argomento di questo capitolo. Questa letteratura si distanzia tanto dal modello francese presente al nord quanto dai temi laici della scuola siciliana
\end{itemize}

La letteratura religiosa viene per molto tempo considerata una letteratura di serie B, questa concezione ci è testimoniata da Dante in due contesti diversi: nel de vulgari eloquentia Dante sostiene che tutti coloro che fanno letteratura si ispirano alla scuola siciliana e, nel paradiso quando Dante incontra San Francesco lo loda per il suo ruolo in ambito religioso senza mai fare accenno alla sua letteratura.
La fama di questa corrente, come vedremo meglio analizzando il testo di Francesco arriva grazie al romanticismo dell'800 che proprio nella figura di Francesco rivede l'esaltazione della natura.

La poesia umbra si distanzia fortemente da quella vista in precedenza, se infatti i poeti della scuola siciliana, soprattutto a causa dell'assolutismo politico potevano trattare solamente tematiche edonistiche e di piacere, la poesia di questa corrente ha un valore maggiormente \textbf{utile} e pragmatico, le poesie sono inserite nei \textit{laudari} con il fine di far pregare i membri di un certo monastero, usano perciò il verso e la metrica non tanto con un fine stilistico ma piuttosto per agevolare la memorizzazione della preghiera.

\subsection{Gli autori}
Anche di questa letteratura abbiamo una scarsa selezione di autori, di cui sappiamo relativamente poco. E' interessante notare come entrambi rappresentino due modelli di letteratura diversissimi se non antitetici tra di loro.
\subsubsection{Francesco d'Assisi}
\subsubsection{Jacopone da Todi}


\end{document}